% !TEX program = xelatex
\documentclass[UTF8, a4paper, 11pt]{ctexart}

% Page layout
\usepackage{geometry}
\geometry{left=2.5cm, right=2.5cm, top=2.5cm, bottom=2.5cm}
\usepackage{fancyhdr}
\usepackage{lastpage}

% Math and relational algebra symbols
\usepackage{amsmath}
\usepackage{amssymb}
\usepackage{mathtools}

% Tables, lists, colors, and links
\usepackage{booktabs}
\usepackage{array}
\usepackage{enumitem}
\usepackage{xcolor}
\usepackage[hidelinks]{hyperref}

% Optional: code-like formatting for SQL fragments or relation schemas
\usepackage{listings}
\lstset{
  basicstyle=\ttfamily\small,
  columns=fullflexible,
  keepspaces=true,
  frame=single,
  breaklines=true
}

% Header and footer
\pagestyle{fancy}
\setlength{\headheight}{13.6pt}
\fancyhf{}
\lhead{数据库导论}
\chead{Homework 2}
\rhead{关系代数}
\cfoot{第 \thepage\ 页，共 \pageref{LastPage} 页}

% Common relational algebra commands
\newcommand{\ra}{\ensuremath{\mathcal{RA}}}
\newcommand{\rel}[1]{\ensuremath{\mathsf{#1}}}
\newcommand{\attr}[1]{\ensuremath{\mathit{#1}}}
\newcommand{\schema}[2]{\ensuremath{\rel{#1}(#2)}}
\newcommand{\pk}[1]{\underline{#1}}
\newcommand{\fk}[1]{#1^{\ast}}

\newcommand{\select}[2]{\ensuremath{\sigma_{#1}\!\left(#2\right)}}
\newcommand{\project}[2]{\ensuremath{\pi_{#1}\!\left(#2\right)}}
\newcommand{\rename}[2]{\ensuremath{\rho_{#1}\!\left(#2\right)}}
\newcommand{\join}[3]{\ensuremath{#2 \mathbin{\bowtie_{#1}} #3}}
\newcommand{\natjoin}[2]{\ensuremath{#1 \mathbin{\bowtie} #2}}
\newcommand{\semijoin}[2]{\ensuremath{#1 \mathbin{\ltimes} #2}}
\newcommand{\antijoin}[2]{\ensuremath{#1 \mathbin{\triangleright} #2}}
\newcommand{\division}[2]{\ensuremath{#1 \div #2}}
\newcommand{\assign}[2]{\ensuremath{#1 \leftarrow #2}}

% Problem environment
\newcounter{problem}
\newenvironment{problem}[1][]{
  \refstepcounter{problem}
  \section*{题目 \theproblem\if\relax\detokenize{#1}\relax\else：#1\fi}
  \addcontentsline{toc}{section}{题目 \theproblem}
}{}

\title{\textbf{《数据库概论》第二次课后作业}}
\author{姓名：李云浩 \quad 学号：241880324}
\date{\today}

\begin{document}

\maketitle

\section*{关系模型}

本次作业使用的关系模型如下。带下划线的属性表示主键，带星号的属性表示外键。

\[\begin{aligned}
  \schema{E: employee}{\pk{person\_ name}, street, city} \\
  \schema{W: works}{\pk{person\_ name}, \pk{company\_ name}, salary} \\
  \schema{C: company}{\pk{company\_ name}, city} \\
  \schema{M: manages}{\pk{person\_ name}, manager\_ name}
\end{aligned}\]

% \section*{记号说明}
% 本文使用的关系代数记号如下：
% \[
% \begin{array}{ll}
%   \select{\theta}{R} & \text{选择：从 } R \text{ 中选出满足条件 } \theta \text{ 的元组} \\
%   \project{A_1,\ldots,A_k}{R} & \text{投影：保留属性 } A_1,\ldots,A_k \\
%   \natjoin{R}{S} & \text{自然连接} \\
%   \join{\theta}{R}{S} & \text{条件连接} \\
%   R - S & \text{差运算} \\
%   R \cup S,\ R \cap S & \text{并、交运算} \\
%   \division{R}{S} & \text{除法运算}
% \end{array}
% \]

\begin{problem}
找出没有上级经理的员工信息。

\[
    \natjoin{(\project{person\_ name}{E} - \project{person\_ name}{M})}{E}
\]
\end{problem}

\begin{problem}
找出与其经理居住在同一城市同一街道的所有员工姓名。

令 \( R = \join{manager\_name = E.person\_name}{\project{manager\_ name, person\_name}{M}}{E}, \\ \quad S = \natjoin{\project{person\_ name}{M}}{E}, 
\quad F = (R.street = S.street) \land (R.city = S.city)\)

\[
    \project{S.person\_ name}{\select{F}{\join{R.person\_ name = S.person\_ name}{R}{S}}}
\]
\end{problem}

\begin{problem}
找出不在 “First Bank Corp." 工作的所有员工的姓名。

\[
  \project{person\_name}{E} - \project{person\_ name}{\select{company\_ name = "First \ Bank \ Corp"}{W}}
\]
\end{problem}

\begin{problem}
找出比 "Small Bank Corp." 的所有员工收入都高的所有员工姓名。  

令 \(R = \select{company\_ name = "Small \ Bank \ Corp"}{W}\)

\[
  \project{person\_ name}{W - R} - \project{person\_name}{\select{R.salary \geq (W-R).salary}{R \times (W - R)}}
\]

\end{problem}

\begin{problem}
找出管理超过两名员工（含两位）的经理姓名。

令 \(M_1 := M, M_2 := M\) 

\(F = (M_1.manager\_name = M_2.manager\_name) \land (M_1.person\_name \neq M_2.person\_name)\)

\[
\project{M_1.manager\_name}{\select{F}{\join{M_1.manager\_name = M_2.manager\_name}{M_1}{M_2}}}
\]


\end{problem}

\begin{problem}
找出在南京所有公司工作过的员工姓名。

\[
{\division{\project{person\_name, company\_name}{W}}{\project{company\_name}{\select{city = "NanJing"}{C}}}}
\]


\end{problem}

\end{document}
